\documentclass[12pt]{article}
 
\usepackage[margin=1.2in]{geometry}
\usepackage{amsmath, amssymb, amsthm}
\usepackage{microtype}
\usepackage{parskip}
\usepackage{titling}
 
% ── Title styling ──────────────────────────────────────────────────────────────
\pretitle{\begin{center}\large\scshape}
\posttitle{\par\end{center}}
\preauthor{}\author{}\postauthor{}
\predate{}\date{}\postdate{}
 
% ── Theorem environments ───────────────────────────────────────────────────────
\newtheoremstyle{problem}
  {12pt}{12pt}{\normalfont}{0pt}{\bfseries}{.}{0.5em}{}
\theoremstyle{problem}
\newtheorem*{problem}{Problem}
 
\newtheoremstyle{solution}
  {12pt}{12pt}{\normalfont}{0pt}{\itshape}{.}{0.5em}{}
\theoremstyle{solution}
\newtheorem*{solution}{Solution}
 
% ── Convenience macros ─────────────────────────────────────────────────────────
\newcommand{\range}{\operatorname{range}}
\newcommand{\F}{\mathbb{F}}
 
\begin{document}
 
\title{Linear Algebra Done Right, 4th Ed. by Sheldon Axler, Problem 1 from Chapter 3C}
\maketitle
\vspace{-1.5em}
\begin{center}
  \rule{0.45\textwidth}{0.4pt}
\end{center}
\vspace{0.5em}
 
% ── Problem ────────────────────────────────────────────────────────────────────
\begin{problem}
Suppose $T \in \mathcal{L}(V, W)$. Show that with respect to each choice of
bases of $V$ and $W$, the matrix of $T$ has at least $\dim \range\, T$ nonzero
entries.
\end{problem}
 
\bigskip
 
% ── Solution ───────────────────────────────────────────────────────────────────
\begin{solution}
Let $v_1, \ldots, v_n$ be a basis of $V$ and $w_1, \ldots, w_m$ be a basis
of $W$, and let $M$ denote the matrix of $T$ with respect to these bases.
By definition, the $j$-th column of $M$ consists of the scalars
$A_{1,j}, \ldots, A_{m,j}$ such that
\[
    T(v_j) \;=\; A_{1,j}\,w_1 + \cdots + A_{m,j}\,w_m.
\]
The span of the columns of $M$ therefore corresponds exactly to the set of
all linear combinations of $T(v_1), \ldots, T(v_n)$, which is $\range\, T$.
Consequently, the number of linearly independent columns of $M$ equals
\[
    \dim \range\, T \;=:\; k.
\]
 
Hence $M$ has at least $k$ linearly independent columns.  Each such column
must be nonzero, since the zero vector cannot belong to any linearly
independent list.  A nonzero column contains at least one nonzero entry, and
entries belong to exactly one column, so these $k$ columns contribute at
least $k$ \emph{distinct} nonzero entries to $M$.
 
Therefore the matrix of $T$ has at least $\dim \range\, T$ nonzero
entries. \hfill$\blacksquare$
\end{solution}
 
\end{document}
 